\subsection{Holomorphic functional calculus in a complete normable real or complex algebra}

Let E be a real topological vector space. The topological vector space \(C \otimes E\), the complexification of E (EVT, II, p.~65), is denoted \(\mathrm{E}_{(\mathbf{C})}\) and E is identified with a real topological vector subspace of \(\mathrm{E}_{(\mathbf{C})}\) by the map \(x \mapsto 1 \otimes x\) .

PROPOSITION 18. --- The complex topological vector space \(E_{(\mathbf{C})}\) is normable (resp. complete) if and only if E is normable (resp. complete).

The underlying real topological vector space of \(\mathrm{E}_{(\mathbf{C})}\) is isomorphic to \(E \times E\) . Thus \(\mathrm{E}_{(\mathbf{C})}\) is complete if and only if E is complete, and E is normable if \(\mathrm{E}_{(\mathbf{C})}\) is.

Suppose conversely that E is normable. Let p be a norm defining the topology of E and B the unit ball of p. There exists a closed balanced neighborhood V of 0 in \(\mathrm{E}_{(\mathbf{C})}\) contained in \(B+iB\) (EVT, II, p.~66). The sets \(\lambda V\) , where \(\lambda\) describes \(R_{+}^{*}\) , therefore form a fundamental system of neighborhoods of 0 in \(\mathrm{E}_{(\mathbf{C})}\) . The gauge of V is a norm on \(\mathrm{E}_{(\mathbf{C})}\) that defines the topology of \(\mathrm{E}_{(\mathbf{C})}\) , hence \(\mathrm{E}_{(\mathbf{C})}\) is normable.

Remark. --- Let E and F be normable topological vector spaces over K. The vector space \(\mathcal{L}(\mathrm{E};\mathrm{F})\) of continuous linear maps from E into F, equipped with the topology of bounded convergence, is a normable topological vector space (EVT, III, p.~14).

Let E and F be normable topological vector spaces over R. The C-linear map \(\varphi\colon\mathcal{L}(\mathrm{E};\mathrm{F})_{(\mathbf{C})}\to\mathcal{L}(\mathrm{E}_{(\mathbf{C})};\mathrm{F}_{(\mathbf{C})})\) defined by \(\varphi(\lambda\otimes u)=\lambda u_{(\mathbf{C})}\) is an isomorphism of complex topological vector spaces. In particular, the dual of \(\mathrm{E}_{(\mathbf{C})}\) is identified with the complexification of the dual of E, and the normed algebra \(\mathcal{L}(\mathrm{E}_{(\mathbf{C})})\) with the complexification of the normed algebra \(\mathcal{L}(\mathrm{E})\) .

Let S be a compact subset of C stable under complex conjugation. Consider the C-algebra \(\mathcal{O}(S)\) of germs of complex-valued holomorphic functions in a neighborhood of S, endowed with the structure of a complex locally convex space defined in No. 1 of I, p.~49. If U is an open neighborhood of S in C, and \(h: U \to C\) be a holomorphic function, the image V of U under complex conjugation is an open neighborhood of S in C and \(h^{*}: w \mapsto \overline{h(\overline{w})}\) is a holomorphic function on V. Passing to the inductive limit yields a continuous involution \(f\mapsto f^{*}\) on the algebra \(\mathcal{O}(\mathrm{S})\) . In particular, we have:

\[
(f + g) ^ {*} = f ^ {*} + g ^ {*} \quad (f g) ^ {*} = f ^ {*} g ^ {*} \quad (\lambda f) ^ {*} = \bar {\lambda} f ^ {*}
\]

for \(f,g\) in \(\mathcal{O}(\mathrm{S})\) and \(\lambda\) in \(\mathbf{C}\) .

We denote by \(\mathcal{O}_{\mathbf{R}}(\mathrm{S})\) the set of germs \(f \in \mathcal{O}(\mathrm{S})\) such that \(f = f^{*}\) . It is a closed full sub- \(\mathbf{R}\) -algebra of \(\mathcal{O}(\mathrm{S})\) .

PROPOSITION 19. --- Let z denote the germ in \(\mathcal{O}(\mathrm{S})\) of the identity mapping of \(\mathbf{C}\) . Then \(\mathcal{O}_{\mathbf{R}}(\mathrm{S})\) is the smallest closed full sub- \(\mathbf{R}\) -algebra of \(\mathcal{O}(\mathrm{S})\) containing \(z\) .

We have \(z^{*} = z\) , hence z belongs to \(\mathcal{O}_{\mathbf{R}}(\mathrm{S})\). Let B be a closed full sub-R-algebra of \(\mathcal{O}(\mathrm{S})\) containing z. The map \(f \mapsto f + f^{*}\) of \(\mathcal{O}(\mathrm{S})\) into \(\mathcal{O}_{\mathbf{R}}(\mathrm{S})\) is continuous and surjective, and the set of germs of rational holomorphic functions in a neighborhood of S is dense in \(\mathcal{O}(\mathrm{S})\) (Theorem 3 of I, p.~69). To show that B contains \(\mathcal{O}_{\mathbf{R}}(\mathrm{S})\) it therefore suffices to show that if f is the germ of such a rational function, then we have \(f + f^{*} \in B\) .

There exist polynomials P and Q in \(\mathbf{C}[X]\) such that Q does not vanish at any point of S and such that we have \(f = \frac{\mathrm{P}(z)}{\mathrm{Q}(z)}\) . Denote by \(P^{*}\) and \(Q^{*}\) the polynomials obtained by replacing the coefficients of P and Q by their conjugates. We then have \(\mathrm{P}(z)^{*} = \mathrm{P}^{*}(z)\) and \(\mathrm{Q}(z)^{*} = \mathrm{Q}^{*}(z)\). Since S is stable under complex conjugation, the polynomial \(Q^{*}\) vanishes at no point of S. The germs \(\mathrm{Q}^{*}(z)\) and \((\mathrm{QQ}^{*})(z)\) are therefore invertible in \(\mathcal{O}(\mathrm{S})\) , and

\[
f + f ^ {*} = \frac {\mathrm{P} (z)}{\mathrm{Q} (z)} + \frac {\mathrm{P} ^ {*} (z)}{\mathrm{Q} ^ {*} (z)} = \frac {(\mathrm{PQ} ^ {*} + \mathrm{P} ^ {*} \mathrm{Q}) (z)}{(\mathrm{QQ} ^ {*}) (z)}.
\]

Since the polynomials PQ* + P*Q and QQ* have real coefficients and B is a full R-subalgebra of \(\mathcal{O}(\mathrm{S})\) containing \(z\) , the element \(f + f^{*}\) belongs to B. This concludes the proof of the proposition.

Let A be a complete normable unital algebra over R. Let x be an element of A. The spectrum of the element \(1 \otimes x\) of the algebra \(\mathrm{A}_{(\mathbf{C})}\) is called the complex spectrum of x, and is denoted \(\mathrm{Sp}_{\mathrm{A}_{(\mathbf{C})}}(x)\) . Its intersection with the set R is none other than the spectrum \(\mathrm{Sp}_{\mathrm{A}}(x)\) of x relative to A, which is sometimes called the real spectrum of x. The complex spectrum \(\mathrm{Sp}_{\mathrm{A}_{(\mathbf{C})}}(x)\) is a compact subset of C, stable under complex conjugation. It is nonempty when the algebra A is not reduced to 0.

Let \(x\) be an element of A. The spectral radius of \(1 \otimes x \in \mathrm{A}_{(\mathbf{C})}\) is equal to the spectral radius \(\varrho(x)\) of \(x\) . It is the smallest real number \(r \geqslant 0\) such that \(|\lambda| \leqslant r\) for all \(\lambda \in \mathrm{Sp}_{\mathrm{A}_{(\mathbf{C})}}(x)\) . We have

\[
\varrho (x) = \lim _ {n \to + \infty} \| x ^ {n} \| ^ {1 / n} = \inf _ {n > 0} \| x ^ {n} \| ^ {1 / n}
\]

for every norm on A that defines the topology of A. Indeed, one may assume that the norm on A is the restriction of a norm on \(\mathrm{A}_{(\mathbf{C})}\) that defines the topology of \(\mathrm{A}_{(\mathbf{C})}\) and apply prop. 1 of I, p.~20.

We denote by \(u \mapsto \overline{u}\) the endomorphism of the R-algebra \(A_{(C)}\) which maps \(\lambda \otimes a\) to \(\overline{\lambda} \otimes a\) . It is continuous.

Lemma. --- For every \(f \in \mathcal{O}(\mathrm{Sp}_{\mathrm{A}_{(\mathbf{C})}}(x))\) , one has \(f^{*}(1 \otimes x) = \overline{f(1 \otimes x)}\) .

The maps \(f \mapsto f(1 \otimes x)\) and \(f \mapsto \overline{f^*(1 \otimes x)}\) are continuous unital homomorphisms of \(\mathbf{C}\) -algebras from \(\mathcal{O}(\mathrm{Sp}(x))\) into \(\mathrm{A}_{(\mathbf{C})}\) which map \(z\) to \(1 \otimes x\); therefore they are equal (I, p.~74, th. 5).

PROPOSITION 20. --- For every \(f \in \mathcal{O}_{\mathbf{R}}(\mathrm{Sp}_{\mathrm{A}_{(\mathbf{C})}}(x))\) , there exists a unique element \(f(x)\) of A such that \(f(1 \otimes x) = 1 \otimes f(x)\) in \(\mathrm{A}_{(\mathbf{C})}\) . The map \(f \mapsto f(x)\) of \(\mathcal{O}_{\mathbf{R}}(\mathrm{Sp}_{\mathrm{A}_{(\mathbf{C})}}(x))\) in A is the unique continuous unital homomorphism of R-algebras that maps the germ in \(\mathcal{O}_{\mathbf{R}}(\mathrm{Sp}_{\mathrm{A}_{(\mathbf{C})}}(x))\) of the identity map of C to \(x\).

We denote \(S = \text{Sp}_{\text{A}(\mathbf{C})}(x)\) . By the above lemma, for every germ \(f \in \mathcal{O}_{\mathbf{R}}(\text{Sp}(x))\) , one has \(f(1 \otimes x) = \overline{f(1 \otimes x)}\). The first assertion follows from this. Let us denote by \(z\) the germ in \(\mathcal{O}_{\mathbf{R}}(\text{S})\) of the identity mapping of \(\mathbf{C}\) . The map \(f \mapsto f(x)\) is a continuous unital homomorphism of the \(\mathbf{R}\) -algebra \(\mathcal{O}_{\mathbf{R}}(\text{Sp}(x))\) into \(A\) , which maps \(z\) to \(x\). This is the unique one by Prop. 19, since any morphism with these properties is uniquely determined on every sub- \(\mathbf{R}\) -algebra of \(\mathcal{O}(\text{S})\) containing \(z\) .

Let \(f \in \mathcal{O}_{\mathbf{R}}(\mathrm{Sp}_{\mathrm{A}(\mathbf{C})}(x))\) . The element \(f(x)\) belongs to every closed full subalgebra of A containing \(x\) (prop. 19), hence belongs to the bicommutant of \(x\) in A. The complex spectrum of \(f(x)\) is equal to \(f(\mathrm{Sp}(x))\) (I, p.~75, prop. 8). For all \(g \in \mathcal{O}_{\mathbf{R}}(f(\mathrm{Sp}_{\mathrm{A}(\mathbf{C})}(x)))\) , one has \(g \circ f \in \mathcal{O}_{\mathbf{R}}(\mathrm{Sp}_{\mathrm{A}(\mathbf{C})}(x))\) and (loc. cit.) \(g \circ f(x) = g(f(x))\) .

Let U be an open subset of C, stable under complex conjugation. The set \(\Omega\) of elements x of A whose complex spectrum is contained in U is open in A (I, p.~76, prop. 10). Let f be a holomorphic function on U such that \(f^{*} = f\) . The map \(x \mapsto f(x)\) of \(\Omega\) in A is analytic (loc. cit.).

Let A, B be complete normable unital associative algebras over R and \(\varphi\colon A\to B\) a continuous unital algebra morphism. Let \(x\in A\) . The complex spectrum of \(\varphi(x)\) is contained in that of x and, for every \(f\in\mathcal{O}_{\mathbf{R}}(\mathrm{Sp}_{\mathrm{A}(\mathbf{C})}(x))\) , one has \(f(\varphi(x))=\varphi(f(x))\) . This follows immediately from the analogous statement in the complex case (I, p.~75, prop. 8).

